\paragraph{\trMatchingCycle{}。}\label{sub:matching-cycles-fibration}
现在我们来更仔细地考察前面那个局势的一个重要特殊情形，即 $L=\Sigma_{\mu,I}$ 是与 $S$ 中某条\trMatchingPath{}$\mu$ 相伴的\trMatchingCycle，而 $\mu$ 相对于诸 $\gamma_k$ 的位置如\cref{fig:matching-triangle}所示。给 $\widetilde{\Delta}_k$、$\widetilde{\Delta}_{k+1}$ 配备\trEquivariant\trBraneStructure，使得 $V_k$、$V_{k+1}$ 上诱导的\trBraneStructure{}成为\trIsotopy{}的（这不过是把为\cref{thm:graded-twist-isotopy}固定诸\trGrading{}的方式重新表述了一遍）。类似地，我们为 $L$ 挑选一个\trGrading，并给 $\widetilde{L}$ 配备与之相应的\trEquivariant\trBraneStructure。
\begin{lemma}\label{thm:matching-delta-hf-sphere}
$HF^{*}(\widetilde{\Delta}^{\#}_k,\widetilde{\Delta}^{\#}_{k+1})\cong H^{*}(S^{n};\KK)$，并且其上的 $\ZZ/2$-作用是平凡的。
\end{lemma}
\begin{proof}
这不过是 \eqref{eq:cf-delta-vanishing-iso} 的一个特殊情形：$HF^{*}(\widetilde{\Delta}^{\#}_k,\widetilde{\Delta}^{\#}_{k+1})\cong HF^{*}(V^{\#}_k,V^{\#}_{k+1})$，且 $\ZZ/2$-作用是平凡的；而由于这两个\trBrane{}是\trIsotopy{}的，后一个群\trIsomorphic{}于 $H^{*}(V_k;\KK)$。
\end{proof}
\begin{lemma}\label{thm:matching-l-delta-hf}
诸群 $HF^{*}(\widetilde{L}^{\#}_{\pm},\widetilde{\Delta}^{\#}_k)$ 都是一维的，并且集中于同一个\trDegree。
\end{lemma}
\begin{proof}
只要认真思考过\trMatchingCycle{}的构造，这一点就应当是清楚的。用\cref{thm:transport-isotopy-realization}，可以找到一个\trConstantAlongFibres{}的\trDeformation，它从给定的\trFibration{}结构出发，最终到达这样一个\trLefschetzFibration：对后者而言，\trMatchingCycle{}可以用朴素的方式来定义，而且其中的\trIsotopy{}$I$ 是常值的。我们可以要求这个\trDeformation{}集中在 $\mu$ 自身的\trPreimage{}附近，从而它\trLift{}为 $\widetilde{E}$ 的一个同类型的\trDeformation。在这个\trDeformation{}后的结构中工作，就会发现所考虑的那些\trLagrange\trSubmanifold{}的类似物——我们把它们称作 $\widetilde{L}^*_{\pm}$ 与 $\widetilde{\Delta}^*_k$——\trTransverse{}相交，并且交于单独一点（该点是 $E_{\gamma_k(1)}$ 中的\trCriticalPoint{}在 $\widetilde{E}$ 中的诸多\trPreimage{}之一）。现在由\cref{thm:moser-deformation-embedding}，\trDeformation{}后的\trDoubleBranchedCover{}附带一个到原来的 $\widetilde{E}$ 之中的\trConformallySymplectic{}\trEmbedding，它在相差\trIsotopy{}的意义下把 $\widetilde{L}^*_{\pm}$ 映到 $\widetilde{L}_{\pm}$、把 $\widetilde{\Delta}^*_k$ 映到 $\widetilde{\Delta}_k$。第一种情形是显然的，因为这样的\trIsotopy{}本来就是\trMatchingCycle{}构造的一部分；第二种情形则可以通过考察在整个\trDeformation{}过程中连续产生的\trVanishingCycle{}族来得到。结果便是：存在我们这些\trLagrange\trSubmanifold{}的一个\trIsotopy，经过它之后它们\trTransverse{}相交于单独一点，这当然蕴含\trFloerCohomology{}的一维性。要看出非平凡的部分位于同一个\trDegree，只需使用 $\ZZ/2$-作用 $\epsilon$ 即可，按定义，它交换两个\trFloer{}群。
\end{proof}
为简单起见，假定 $L$ 的\trGrading{}已经选得使诸群 $HF^{*}(\widetilde{L}^{\#}_{\pm},\widetilde{\Delta}^{\#}_k)$ 都集中在\trDegree{}$1$。
\begin{lemma}\label{thm:matching-exact-triangle}
在 $H(\widetilde{\scr{A}})$ 中有如下形状的一个\trExactTriangle
\begin{equation}\label{eq:matching-exact-triangle-diagram}
\begin{tikzcd}
\widetilde{L}^{\#} \ar[rr,"{[1]}"] && \widetilde{\Delta}^{\#}_k \ar[dl] \\
& \widetilde{\Delta}^{\#}_{k+1}. \ar[ul] &
\end{tikzcd}
\end{equation}
\end{lemma}
\begin{proof}
由\cref{thm:graded-twist-isotopy}与\cref{thm:twist-equals-tau}我们知道，在 $D(\widetilde{\scr{A}})$ 中，
\begin{equation}\label{eq:matching-twist-iso}
\widetilde{\Delta}^{\#}_{k+1}\cong T_{\widetilde{L}^{\#}_{-}}T_{\widetilde{L}^{\#}_{+}}\bigl(\widetilde{\Delta}^{\#}_{k}\bigr).
\end{equation}
此外，诸 $\widetilde{L}_{\pm}$ 互不相交，因此它们之间没有\trMorphism，这意味着
\begin{equation*}
T_{\widetilde{L}^{\#}_{-}}T_{\widetilde{L}^{\#}_{+}}\bigl(\widetilde{\Delta}^{\#}_{k}\bigr)
\cong T_{\widetilde{L}^{\#}_{-},\widetilde{L}^{\#}_{+}}\bigl(\widetilde{\Delta}^{\#}_{k}\bigr),
\end{equation*}
其中右边正如\cref{thm:multiple-twisting-objects}、\cref{thm:twist-pair-homology-triangle}中所定义的那样。记住这一点，\eqref{eq:matching-twist-iso}便在 $H(Tw\widetilde{\scr{A}})$ 中导出如下形状的一个\trExactTriangle
\begin{equation*}
\begin{tikzcd}
HF^{*}(\widetilde{L}_{-},\widetilde{\Delta}^{\#}_k)\otimes\widetilde{L}_{-}\oplus
HF^{*}(\widetilde{L}_{+},\widetilde{\Delta}^{\#}_k)\otimes\widetilde{L}_{+}
\ar[rr] && \widetilde{\Delta}^{\#}_k \ar[dl] \\
& \widetilde{\Delta}^{\#}_{k+1}. \ar[ul,"{[1]}"'] &
\end{tikzcd}
\end{equation*}
由\cref{thm:matching-l-delta-hf}，左上角的那个\trVertex{}是 $S^{-1}(\widetilde{L}^{\#}_{-})\oplus S^{-1}(\widetilde{L}^{\#}_{+})\cong S^{-1}(\widetilde{L}^{\#})$。现在把这个三角形旋转（\cref{thm:rotated-exact-triangle}）。
\end{proof}
\begin{proposition}\label{thm:matching-l-cone}
在 $D\scr{F}(\pi)$ 中，$\widetilde{L}^{\#}$\trIsomorphic{}于 $\mathrm{Cone}(\widetilde{\Delta}^{\#}_k\to\widetilde{\Delta}^{\#}_{k+1})$，其中这个锥可以在任何一个非平凡的零次\trMorphism{}上作成。
\end{proposition}
\begin{proof}
由 \eqref{eq:matching-exact-triangle-diagram}，我们得到 $D(\widetilde{\scr{A}})$ 中的一个\trIsomorphism{}$\widetilde{L}^{\#}\cong\mathrm{Cone}(\widetilde{\Delta}^{\#}_k\to\widetilde{\Delta}^{\#}_{k+1})$。在\trIsomorphic{}的意义下，任何\trMorphism{}的锥对于把该\trMorphism{}乘以 $\KK^{\times}$ 中某个元素的操作是不变的。在我们目前的情形，\trMorphism{}群是一维的，于是剩下有待证明的只是：所用的那个\trMorphism{}非零。倘若它是零，两个\trDirectSummand{}$\widetilde{L}^{\#}_{\pm}$ 就不得不分别\trIsomorphic{}于 $\widetilde{\Delta}^{\#}_k$、$\widetilde{\Delta}^{\#}_{k+1}$，但这立即与我们先前的\trFloerCohomology{}计算相矛盾（例如\cref{thm:matching-l-delta-hf}）。最后，由于 $\widetilde{L}^{\#}$ 是自由的，我们可以应用\cref{thm:equivariant-cone-isomorphism}来得到 $D(\widetilde{\scr{A}}^{\mathrm{inv}})=D\scr{F}(\pi)$ 中的一个\trIsomorphism。
\end{proof}
\begin{remark}\label{thm:matching-alternative-approach}
另一种做法是把当前这个局势严格当作\cref{thm:type-u-in-generated-subcategory}的特殊情形来处理。为简单起见，设 $k=1$。这时需要证明（通过考察一个适当的局部模型）：乘积
\begin{equation*}
HF^{*}(\widetilde{\Delta}^{\#}_2,\widetilde{L}^{\#}_{+})\otimes HF^{n}(\widetilde{\Delta}^{\#}_1,\widetilde{\Delta}^{\#}_2)
\longrightarrow HF^{*+n}(\widetilde{\Delta}^{\#}_1,\widetilde{L}^{\#}_{+})
\end{equation*}
是一个\trIsomorphism（所涉及的每个空间实际上都是一维的）。由此，再加上 $HF^{*}(\widetilde{\Delta}^{\#}_j,\widetilde{L}^{\#}_{+})=0$（$j>2$）这一事实，便可推知：\cref{thm:equivariant-split-generation}之证明中的那个\trObject{}$U$ 在当前这个特殊情形下化为 $\mathrm{Cone}(\widetilde{\Delta}^{\#}_1\to\widetilde{\Delta}^{\#}_2)$。
\end{remark}
